\documentclass[10pt,a4paper]{amsart}
\usepackage[margin=19mm]{geometry}
\usepackage{amsmath,amssymb,mathtools}
\usepackage[scr=rsfs,cal=euler]{mathalfa}
\usepackage{xfrac}
\usepackage[unicode,hidelinks,bookmarks=false]{hyperref}
\newcommand{\ie}{i.e.\ }
\newcommand{\cf}{cf.\ }
\newcommand{\resp}{respectively\ }
\newcommand{\Subsection}{\subsection}
\providecommand{\nobreakdash}{\nobreak\hbox{-}}
\setlength{\emergencystretch}{2em}
\multlinegap0pt
\usepackage{bm}
\usepackage{bbm}
\usepackage{dsfont}
\usepackage{xspace}
\usepackage{cancel}
\usepackage{mathtools}
\usepackage{amsthm}
\makeatletter
\@ifundefined{lemma}{\newtheorem{lemma}{Lemma}}{}
\@ifundefined{proposition}{\newtheorem{proposition}[lemma]{Proposition}}{}
\makeatother
\makeatletter
\expandafter\ifx\csname proof\endcsname\relax
\newenvironment{proof}{{\it Proof.}~~}{\hfill$\square$}
\fi
\makeatother
\title[Exponential convergence of multiagent systems with lack of c]{Exponential convergence of multiagent systems with lack of connection}
\author{Fabio Ancona, Mohamed Bentaibi, Francesco Rossi}
\date{Source: arXiv:2510.11334v1 (CC BY 4.0)}
\begin{document}
\maketitle

\noindent\emph{Source and licence.} Exponential convergence of multiagent systems with lack of connection. Version arXiv:2510.11334v1 (CC BY 4.0), distributed under the Creative Commons Attribution 4.0 International licence, as recorded in the arXiv licence metadata for this exact version and shown on the abstract page https://arxiv.org/abs/2510.11334v1. Full rights notice: licenses/arxiv-2510.11334-CC-BY-4.0.txt.\par\smallskip

\noindent\textit{Editorial scope.} Counted unit: the Proposition whose source label is \texttt{p-1step} in the article (source0.tex lines 453--459, source label \texttt{p-1step}), delivered together with the article's own complete original proof of it (source0.tex lines 460--525, one \texttt{proof} environment of 66 lines). No mathematics is written, repaired or re-derived: statement and proof are the article's own text line for line. Required context retained uncounted: e-ODEmain (lines 401--404); p-barrier (lines 416--419); e-stimaf (lines 422--436). Every definition, display and label of those lines is retained unchanged. Omitted from this package and not redistributed: 1--400, 405--415, 420--421, 437--452, 526--1156. Nothing inside a retained excerpt is omitted or altered: every retained body is delivered exactly as the article's lines are written, so no omission receipt of any kind accompanies this package. Citations: the retained excerpts carry no citation command at all, so no bibliography entry of the article is transcribed and no reference is redistributed. Reference adaptation: none was needed and none was made; every reference inside the delivered unit resolves inside it, so no unresolved reference or citation is produced. Printed slips kept literal: no printed word, symbol, hypothesis or number of the counted statement or of its proof was altered. Layout and class: the article's own class is \texttt{article} with packages including graphicx, color, amsfonts, amsmath, mathrsfs, hyperref, lipsum, epstopdf, algorithmic, cancel, fullpage, physics, tikz, mathdots; the wrapper is the lane's pinned \texttt{amsart} editorial wrapper at 10pt on A4, which restates the theorem-like environments the retained text needs and adds the article's own macro definitions that the retained text needs (none), with extra wrapper packages (bm, bbm, dsfont, xspace, cancel, mathtools). The delivered numbering is the unit's own. Assets: no retained span contains a figure environment, a graphics-inclusion command, a PDF-page inclusion, a table or a TikZ picture; no image file is copied into this package, the package declares no asset dependency and no image file is redistributed with it. Licence: version arXiv:2510.11334v1 of this article is distributed under the Creative Commons Attribution 4.0 International licence, as recorded in the arXiv licence metadata for this exact version and shown on the abstract page https://arxiv.org/abs/2510.11334v1. A full rights notice is delivered at licenses/arxiv-2510.11334-CC-BY-4.0.txt. Characters: every retained excerpt is pure ASCII, so the delivered engine, XeLaTeX with its default Latin Modern text font, reports no missing character.
\begin{equation}
\label{e-ODEmain}
\dot x_i(t)=\frac{1}{N}\sum_{j=1}^N M_{ij}(t) (x_j-x_i)\mbox{~~with~~}M_{ij}:[0,+\infty)\mapsto [0,1],
 \end{equation}

\begin{proposition}[Left barrier]\label{p-barrier} Define $$\Psi_L(\bar x,\alpha,t):=\alpha+\exp\left(-\tfrac{N-1}N t\right)(\bar x-\alpha).$$
Let $x(t)$ be a solution of \eqref{e-ODEmain} with $x_i(0)\geq \alpha$ for all $i\in\{1,\ldots,N\}$,
and let $\bar x \geq\alpha$. If there exists an index $I$ and a time $\tau\in[0,T]$ such that $x_I(\tau)\geq \Psi_L(\bar x,\alpha,\tau)$, then for all $t\geq \tau$ it also holds $x_I(t)\geq \Psi_L(\bar x,\alpha,t)$.
\end{proposition}

\begin{eqnarray}
\frac{d}{dt} f(t)&=&\frac1N \sum_{x_j(t) \leq \Psi_L(\bar x,\alpha,t),j\neq I} M_{Ij}(t)(x_j(t)-x_I(t))+\frac1N \sum_{x_j(t)> \Psi_L(\bar x,\alpha,t),j\neq I} M_{Ij}(t)(x_j(t)-x_I(t))+\nonumber \\
&& \tfrac{N-1}N \exp\left(-\tfrac{N-1}N t\right)(\bar x-\alpha)\nonumber \\
%
&\geq& \frac1N \sum_{x_j\leq \Psi_L(\bar x,\alpha,t),j\neq I} M_{Ij}(t)\bigg( \Big(\alpha-\Psi_L(\bar x,\alpha,t) \Big)   + \Big(\Psi_L(\bar x,\alpha,t)-x_I(t) \Big) \bigg) +\nonumber\\
&& \hspace{-5mm} \frac1N \sum_{x_j(t)> \Psi_L(\bar x,\alpha,t),j\neq I} M_{Ij}(t) \Big(\Psi_L(\bar x,\alpha,t)-x_I(t)\Big)  + \tfrac{N-1}N \exp\left(-\tfrac{N-1}N t\right)(\bar x-\alpha)\nonumber\\
&=& \frac1N \sum_{x_j\leq \Psi_L(\bar x,\alpha,t),j\neq I} M_{Ij}(t)\Big(\alpha-\Psi_L(\bar x,\alpha,t) \Big) + \frac1N \sum_{j\neq I} M_{Ij}(t)(\Psi_L(\bar x,\alpha,t)-x_I(t))+ \nonumber \\
%
&& \tfrac{N-1}N \exp\left(-\tfrac{N-1}N t\right)(\bar x-\alpha)\nonumber \\
%
&\geq &\frac1N \sum_{j\neq I} M_{Ij}(t) \Big(-  \exp\left(-\tfrac{N-1}N t\right)  (\bar x-\alpha) \Big)  - \bigg(\frac1N \sum_{j\neq I} M_{Ij}(t) \bigg) f(t) + \tfrac{N-1}N \exp\left(-\tfrac{N-1}N t\right)(\bar x-\alpha)\nonumber\\
&= &\frac1N \sum_{j\neq I} (1-M_{Ij}(t)) \exp\left(-\tfrac{N-1}N t\right)(\bar x-\alpha) 
 - \bigg(\frac1N \sum_{j\neq I} M_{Ij}(t) \bigg) f(t) \geq 0-r(t) f(t).
\label{e-stimaf}
\end{eqnarray}

\begin{proposition}\label{p-1step} Let $T,\mu>0$ be given. Let $x(t)$ be a solution of \eqref{e-ODEmain} with $x_k(0)\geq \alpha$ for all $k\in\{1,\ldots,N\}$. Assume that the $(T,\mu)$-connectivity graph $G(t)$ contains the directed edge $i\to j$. It then holds
\begin{eqnarray}
\label{e-1step}
x_i(T)\geq \alpha+\eta\exp(-2\tfrac{N-1}{N}T)(x_j(0)-\alpha)
\end{eqnarray}
with $\eta=\frac{\mu T}{N+\mu T}$.
\end{proposition}

\begin{proof}
We first define the following barrier: $$\gamma_j(t):=\Psi_L(x_j(0),\alpha,t)=\alpha+\exp(-\tfrac{N-1}{N}t)(x_j(0)-\alpha).$$
As a first case, assume that $x_i(0)\geq \gamma_j(T)$. We rewrite it as $x_i(0)\geq \Psi_L(\gamma_j(T),\alpha,0)$  and apply Proposition \ref{p-barrier}. Remark that we can apply such proposition
to the left barrier $\gamma_j$, since $\gamma_j(T)\geq \alpha$ as a consequence of $x_j(0)\geq\alpha$. It then holds
\begin{eqnarray*}
x_i(T)&\geq& \Psi_L(\gamma_j(T),\alpha,T)\geq \alpha+\exp(-\tfrac{N-1}{N}T)(\alpha+\exp(-\tfrac{N-1}{N}T)(x_j(0)-\alpha)-\alpha)=\\
&&\alpha+\exp(-2\tfrac{N-1}{N}T)(x_j(0)-\alpha).
\end{eqnarray*}
This condition is stronger than the statement, since $1>\eta$.


We then assume that it holds $x_i(0)<\gamma_j(T)$ from now on. We define 
$$l:=\left(1+\frac{\mu T}N\right)^{-1}(\gamma_j(T)-x_i(0)),$$
which is strictly positive by hypothesis. We now define a second barrier function, that is
\begin{eqnarray*}
&&\gamma_i(t):=\Psi_L(\gamma_j(T)-l,\alpha,t)=\alpha+\exp(-\tfrac{N-1}{N}t)(\gamma_j(T)-l-\alpha)=\\
&&\alpha+\exp(-\tfrac{N-1}{N}(t+T))(x_j(0)-\alpha)-\exp(-\tfrac{N-1}{N}t)l.
\end{eqnarray*}

We have two cases: either $x_i(t)\geq \gamma_i(t)$ for some $t\in[0,T]$ or $x_i(t)< \gamma_i(t)$ for all $t\in[0,T]$. In the first case, by the definition of $l$ and recalling that $x_i(0),x_j(0)\geq\alpha$, we first observe the following estimate:
\begin{eqnarray*}
&&\exp(-\tfrac{N-1}{N}T)(x_j(0)-\alpha)\geq \left(1+\frac{\mu T}N\right)^{-1}\exp(-\tfrac{N-1}{N}T)(x_j(0)-\alpha) \geq l.
\end{eqnarray*}
This implies 
$\gamma_j(T)-l\geq \alpha$. We can then apply Proposition~\ref{p-barrier} to the left barrier $\gamma_i$, that ensures $x_i(T)\geq\gamma_i(T)$. This reads as
\begin{eqnarray*}
x_i(T)&\geq& \alpha+\exp(-2\tfrac{N-1}{N}T)(x_j(0)-\alpha)\\
&&-\exp(-\tfrac{N-1}{N}T)\left(1+\frac{\mu T}N\right)^{-1}(\alpha+\exp(-\tfrac{N-1}{N}T)(x_j(0)-\alpha)-x_i(0))\\
&=& \alpha+\exp(-2\tfrac{N-1}{N}T)\eta (x_j(0)-\alpha)+\exp(-\tfrac{N-1}{N}T)\left(1+\frac{\mu T}N\right)^{-1}(x_i(0)-\alpha)\\
&\geq& \alpha+\exp(-2\tfrac{N-1}{N}T)\eta (x_j(0)-\alpha).
\end{eqnarray*}
Here, we used that $x_i(0)-\alpha\geq 0$ by hypothesis on the support, and that $\eta=1-\left(1+\frac{\mu T}N\right)^{-1}$. Also in this case, the statement is proved.

We are then left with the case $x_i(t)< \gamma_i(t)$ for all $t\in[0,T]$. We prove that this case cannot occur, by contradition. We first observe that, due to the fact that left barriers are decreasing functions of time, for all $t\in[0,T]$ it holds both $x_i(t)< \gamma_j(T)-l$ and $x_j(t)\geq \gamma_j(T)$, i.e. $x_j(t)-x_i(t)\geq l$. We now provide an estimate similar to \eqref{e-stimaf} with this added information. We define $f(t):=x_i(t)-\gamma_i(t)$, that is strictly positive, continuous and differentiable for almost every $t\in[0,T]$. We also observe that the following identity holds:
$$\dot \gamma_i(t)=-\tfrac{N-1}{N}\exp(-\tfrac{N-1}{N}t)(\gamma_i(0)-\alpha).$$

 For points of differentiability, it then holds
\begin{eqnarray}
\frac{d}{dt} f(t)&=&\frac1N M_{ij}(t)(x_j(t)-x_i(t))+\frac1N \sum_{x_k(t) \leq \gamma_i(t),k\neq i} M_{ik}(t)(x_k(t)-x_i(t))\nonumber \\
& &+\frac1N \sum_{k\not\in\{ i,j\}, x_k(t)> \gamma_i(t)} M_{ik}(t)(x_k(t)-x_i(t))+\tfrac{N-1}{N}\exp(-\tfrac{N-1}{N}t)(\gamma_i(0)-\alpha)\nonumber \\
&\geq& \frac1N M_{ij}(t)l +\frac1N \sum_{x_k(t) \leq \gamma_i(t),k\neq i} M_{ik}(t)((\alpha-\gamma_i(t))+(\gamma_i(t)-x_i(t)))\nonumber \\
& &+\frac1N \sum_{k\not\in\{i, j\}, x_k(t)> \gamma_i(t)} M_{ik}(t)(\gamma_i(t)-x_i(t))+\tfrac{N-1}{N}\exp(-\tfrac{N-1}{N}t)(\gamma_i(0)-\alpha)\nonumber \\
&\geq& \frac1N M_{ij}(t)l +\frac1N \sum_{x_k(t) \leq \gamma_i(t),k\neq i} M_{ik}(t)(\alpha-\gamma_i(t))+\frac1N \sum_{k\not\in\{i, j\}} M_{ik}(t)(\gamma_i(t)-x_i(t))+\nonumber\\
&&\tfrac{N-1}{N}\exp(-\tfrac{N-1}{N}t)(\gamma_i(0)-\alpha).\nonumber 
\end{eqnarray}
We now observe that it holds $$\alpha-\gamma_i(t)= -\exp(-\tfrac{N-1}{N}t)(\gamma_i(0)-\alpha).$$ We then have 
\begin{eqnarray}
\frac{d}{dt} f(t)&\geq& \frac1N M_{ij}(t)l +\frac1N \sum_{k\neq i} M_{ik}(t)(-\exp(-\tfrac{N-1}{N}t)(\gamma_i(0)-\alpha)) \nonumber\\
&&-\left(\frac1N \sum_{k\not\in\{i, j\}} M_{ik}(t)\right) f(t) +\tfrac{N-1}{N}\exp(-\tfrac{N-1}{N}t)(\gamma_i(0)-\alpha)\nonumber\\
&\geq&\frac1N M_{ij}(t)l + \frac1N \sum_{k\neq i} (1-M_{ik}(t))\exp(-\tfrac{N-1}{N}t)(\gamma_i(0)-\alpha)) -\left(\frac1N \sum_{k\neq i} M_{ik}(t)\right) f(t)\nonumber\\
&\geq& \frac1N M_{ij}(t)l +0- \frac{N-1}N f(t).
\label{e-stimaf2}
\end{eqnarray}
Again, here we used $M_{ik}(t)\in [0,1]$. A direct application of  Gr\"onwall's lemma now implies 
\begin{eqnarray*}
&&f(T)\geq \exp(-\tfrac{N-1}{N}T)\left[ f(0)+\int_0^T\exp(s)\frac1N M_{ij}(s) l\,ds\right]\\
&&\geq \exp(-\tfrac{N-1}{N}T) \left[x_i(0)-\gamma_i(0)+1\cdot \frac l N  \int_0^T
%\blu{\cancel{{\color{black}\frac1N}}} 
M_{ij}(s)\,ds\right].
\end{eqnarray*}
Since $i\to j$ is an arrow of $G(t)$, it holds $\frac1T \int_{0}^T M_{ij}(s)\,ds\geq \mu$. It then holds
\begin{eqnarray*}
x_i(T)-\gamma_i(T)&=&f(T)\geq \exp(-\tfrac{N-1}{N}T) \left(x_i(0)-\left(\gamma_j(T)-l\right)+ \frac{l\, T}N \mu\right)=\\
&&\exp(-\tfrac{N-1}{N}T) \left(x_i(0)-\gamma_j(T)+l\left(1+\frac{\mu\, T}N\right)\right)=0.
\end{eqnarray*}
This raises a contradiction with the condition $x_i(t)< \gamma_i(t)$ for all $t\in[0,T]$.\end{proof}

\end{document}
